\chapter{Contributions}
\label{chap:contributions}

\section*{Introduction}
This chapter explains what was actually built in the project. The main contribution is a
two-phase Unity simulation where reinforcement learning is used for endoscope
navigation and snare capture. I keep the wording simple here on purpose, because the
important part is the system itself, not fancy language.

The work is split into two parts:

\begin{itemize}
	\item Phase~1 learns how to steer the endoscope through a colon-like path while
	forward movement is scripted.
	\item Phase~2 learns how to control the snare and capture a predefined polyp object
	using discrete actions and reward shaping.
\end{itemize}

This setup helped keep the problem manageable. Instead of trying to solve everything at
once, the project separates navigation from capture, which makes training, debugging,
and evaluation much easier.

Another contribution is the way the environment is structured. The agents use only
camera input, the rewards are designed around safety and progress, and the capture task
uses layered gating so the policy cannot just rush into the target without learning the
right sequence.

Finally, the report and code documentation were written so the experiments can be
reproduced more easily. That includes the PPO configuration files, the network
structure, and the main Unity scripts used in each phase.
 

\section{Overview}
\addcontentsline{toc}{section}{Overview}
This section expands the description of the implemented two-phase reinforcement
learning system. The goal is to learn endoscope navigation and an associated
snare-capture behavior in a simulated Unity environment. Important clarification:
the system does \emph{not} perform polyp detection or classification. The agent
policies are vision-only and map tip-camera images to discrete control actions.

\section{System Architecture}
\addcontentsline{toc}{section}{System Architecture}
The environment is implemented in Unity with ML-Agents for training. Key components
include `EndoscopeController` (tip kinematics and collision handling),
`EndoscopeActionDriver` (action mapping), `EndoscopeBurstPusher` (scripted insertion
for Phase~1), `PolypCaptureZone` and `PolypContactSensor` (layered capture logic), and
`SnareSlideController`.

The work is split into two phases to simplify exploration and evaluation:
\begin{itemize}
	\item Phase~1: navigation (steering only, scripted forward push).
	\item Phase~2: snare capture (steering + snare extension/rotation/open-close).
\end{itemize}

\section{Observations and Actions}
\addcontentsline{toc}{section}{Observations and Actions}
Both agents use a single tip-mounted RGB camera rendered at 80$\times$80 pixels. No
external detectors or segmentation masks are provided to the agents; policies learn
directly from images.

Phase~1 action space: one discrete branch with five actions (none, right, down,
left, up) that produce steering inputs. Forward motion is scripted by
`EndoscopeBurstPusher`.

Phase~2 action space: four discrete branches controlling steering (5 actions), snare
slide (3 actions), snare rotation (3 actions), and open/close (3 actions).

\section{Reward Design and Safety}
\addcontentsline{toc}{section}{Reward Design and Safety}
Rewards are shaped to encourage progress while discouraging unsafe contacts. Phase~1
uses a small alive reward, a goal bonus on reaching a trigger, and collision
penalties. Phase~2 uses layered rewards via `PolypCaptureZone` (increasing rewards
for progressing through layers), extension progress bonuses, and contact-based
outcomes (goal sphere = success, large sphere = penalty, colon contact = failure).

The report documents the exact numeric coefficients used during experiments; these
are reproduced in the `config/` YAML files.

\section{Network Architectures and Training}
\addcontentsline{toc}{section}{Network Architectures and Training}
Policies are trained with Proximal Policy Optimization (PPO). The visual encoder is a
small convolutional network that produces a 256-dimensional embedding, followed by
two 256-unit MLP layers and an LSTM memory module (sequence length 64, memory size
128). Policy heads are multi-categorical to match discrete branches. Phase~2 may
use an auxiliary curiosity signal to aid exploration.

Training hyperparameters (batch sizes, buffers, learning rates, schedules) are
stored in `config/endoscope_nav_ppo.yaml` and `config/snare_capture_ppo.yaml`.

\section{Evaluation Protocol}
\addcontentsline{toc}{section}{Evaluation Protocol}
Evaluation uses deterministic checkpoints and fixed seeds. Reported metrics include
success rate, collision rate, time-to-success, and layer progression statistics.
For reproducibility, list the number of seeds, evaluation episode count, and the
policy selection heuristic (best validation checkpoint vs last checkpoint).

\section{Limitations}
\addcontentsline{toc}{section}{Limitations}
This simulation uses Unity rigid-body physics and simplified spherical polyps. The
environment is intended for algorithmic investigation and not for clinical claims.
Where appropriate, the discussion in Chapter~\ref{chap:state_of_the_art} cites the
clinical and perception literature to contextualize results.

\section{Reproducibility}
\addcontentsline{toc}{section}{Reproducibility}
To reproduce experiments install the ML-Agents Python trainers version used during
development and match the Unity editor version. Include `config/` files, the exact
checkpoint hashes, and scripts to run evaluation.

\section{Files of Interest}
\addcontentsline{toc}{section}{Files of Interest}
See `Assets/` for Unity scripts and `config/` for trainer settings. Trainer source
details referenced in the repository's `.venv` folder are helpful for exact network
semantics.
